% !TeX root = ../../../progress-sem1.tex
% 4.x สรุปผลการดำเนินงานเทียบกับแผน (semester 1)
\section{สรุปผลการดำเนินงานเทียบกับแผน}
\thaijustify{
  ตารางที่~\ref{tab:progress-status} สรุปสถานะของงานใน milestone M0 ถึง M3 ซึ่งเป็นแผนของภาคการศึกษาที่ 1 ณ วันที่ 26 พฤศจิกายน 2569 โดยพิจารณาจากซอร์สโค้ดใน branch หลักของ repository
}
\begingroup
\small
\begin{xltabular}{\linewidth}{|>{\RaggedRight\arraybackslash}X|>{\RaggedRight\arraybackslash}p{0.15\linewidth}|>{\RaggedRight\arraybackslash}p{0.36\linewidth}|}
  \caption{สถานะของงานในภาคการศึกษาที่ 1}\label{tab:progress-status}\\
  \hline
  \textbf{งาน} & \textbf{สถานะ} & \textbf{หมายเหตุ} \\ \hline
  \endfirsthead
  \multicolumn{3}{l}{\textit{ตารางที่ \thetable\ (ต่อ)}}\\
  \hline
  \textbf{งาน} & \textbf{สถานะ} & \textbf{หมายเหตุ} \\ \hline
  \endhead
  \multicolumn{3}{|l|}{\textbf{M0: Foundation}} \\ \hline
  ระบุกลุ่มเป้าหมาย ขอบเขต และ MVP & แล้วเสร็จ & บทที่ 1 \\ \hline
  กำหนด data model & แล้วเสร็จ & 34 ตาราง ครอบคลุมทุก milestone (หัวข้อ~\ref{sec:er-diagram}) \\ \hline
  เลือก architecture: auth, LLM provider, prototype rendering, hosting & ดำเนินการบางส่วน & เลือกแล้วทุกส่วนยกเว้นบริการที่ใช้ deploy \\ \hline
  สำรวจและเลือก design system และ UI kit & แล้วเสร็จ & shadcn/ui สำหรับแพลตฟอร์ม และ UI kit shadcn/ui กับ MUI สำหรับ AI \\ \hline
  กำหนด requirements ด้าน PDPA และข้อความขอ consent & แล้วเสร็จ & เก็บความยินยอมพร้อมเวอร์ชันของข้อความ \\ \hline
  ตั้งค่า CI/CD, environments และ event analytics & ดำเนินการบางส่วน & มี CI และสภาพแวดล้อมสำหรับพัฒนาแล้ว ยังไม่มี CD และ event analytics \\ \hline
  การกำหนดดัชนีชี้วัดผล และแผนการดำเนินงาน & แล้วเสร็จ &  \\ \hline
  \multicolumn{3}{|l|}{\textbf{M1: Secure Access}} \\ \hline
  จัดทำ Landing page & แล้วเสร็จ &  \\ \hline
  จัดทำ Register page & ดำเนินการบางส่วน & หน้าจอเสร็จแล้ว backend อยู่ใน branch แยก \\ \hline
  การยืนยันตัวตนผ่านอีเมล & ดำเนินการบางส่วน & อยู่ระหว่างพัฒนาใน branch แยก \\ \hline
  การเข้าสู่ระบบพร้อมระบบจัดการข้อผิดพลาด & ดำเนินการบางส่วน & เข้าสู่ระบบด้วย Google ได้แล้ว อีเมลและรหัสผ่านอยู่ใน branch แยก \\ \hline
  ระบบรีเซ็ตรหัสผ่าน & ดำเนินการบางส่วน & หน้าจอเสร็จแล้ว backend อยู่ใน branch แยก \\ \hline
  ระบบควบคุมการเข้าถึง & แล้วเสร็จ & ควบคุมสิทธิ์ตาม permission ของแต่ละบทบาท \\ \hline
  จัดทำระบบ routing ตามบทบาทของผู้ใช้ & ดำเนินการบางส่วน & backend ตรวจสิทธิ์ตามบทบาทแล้ว หน้าจอของ Client, Developer และ PM อยู่ใน M6 \\ \hline
  ทดสอบระบบยืนยันตัวตนและการควบคุมการเข้าถึง & ดำเนินการบางส่วน & มี unit test และ end-to-end test ของ auth และสิทธิ์ ยังขาดส่วนอีเมลและรหัสผ่าน \\ \hline
  \multicolumn{3}{|l|}{\textbf{M2: Project Setup}} \\ \hline
  ระบบสร้าง workspace & แล้วเสร็จ & รวมการเชิญสมาชิกและรูปประจำ workspace \\ \hline
  ระบบสร้างโปรเจกต์ & แล้วเสร็จ & รวมการกำหนดสมาชิกและบทบาทในโปรเจกต์ \\ \hline
  ฟอร์มกำหนดความต้องการทางการออกแบบ & แล้วเสร็จ &  \\ \hline
  ระบบจัดเก็บความต้องการทางการออกแบบเพื่อใช้เป็น reusable AI context & แล้วเสร็จ & เก็บเวอร์ชันของ design intent ที่ใช้ในแต่ละงาน \\ \hline
  \multicolumn{3}{|l|}{\textbf{M3: AI Design Generation}} \\ \hline
  การสร้าง User story และ Problem statement & แล้วเสร็จ & การให้ AI ร่าง user story มีเฉพาะ backend \\ \hline
  เครื่องมือสร้างคำสั่งที่รวมความต้องการทางการออกแบบและ user story & แล้วเสร็จ &  \\ \hline
  การสร้าง two-alternative UI & แล้วเสร็จ & มีการตรวจและให้โมเดลแก้ไขอัตโนมัติ \\ \hline
  การปรับปรุง UI (Chat-based) & แล้วเสร็จ & เก็บประวัติแต่ละเวอร์ชัน \\ \hline
  การดูตัวอย่างรูปแบบหน้าจอ & แล้วเสร็จ &  \\ \hline
  การอนุมัติและมอบหมายงานให้ผู้เข้าร่วมทดสอบ & ดำเนินการบางส่วน & อนุมัติ variant ได้แล้ว การมอบหมายงานมีเฉพาะ backend \\ \hline
  การเริ่มเซสชันการทดสอบ & ดำเนินการบางส่วน & backend สร้างลิงก์และรหัสได้แล้ว หน้าจอยังไม่มี \\ \hline
  ชุดข้อมูลสำหรับการประเมินคุณภาพผลลัพธ์ของ AI & ยังไม่เริ่ม &  \\ \hline
\end{xltabular}
\endgroup
\thaijustify{
  นอกจากงานตามแผนแล้ว ทีมได้พัฒนาส่วน backend ของ M4 ถึง M6 ไว้ล่วงหน้า ได้แก่ เซสชันการทดสอบพร้อมการเก็บความยินยอม การบันทึกภารกิจ พฤติกรรม และข้อเสนอแนะของผู้เข้าร่วมทดสอบ การวิเคราะห์ข้อเสนอแนะด้วย AI การจัดการประเด็นและบันทึกการตัดสินใจ รวมถึงมุมมองและรายงานสำหรับผู้มีส่วนได้ส่วนเสีย ส่วนหน้าจอของงานเหล่านี้อยู่ในแผนของภาคการศึกษาที่ 2
}
